% !PRINTED 356
\weberpage{356}
%<<CONT
gleich Null für alle der Gleichung $f(x)=0$ genügenden Werthe von $(x)$, und folglich muss sie durch $f(x)$ theilbar sein. Wir können also, wenn wir mit $\Phi(x,a,b)$ eine Form $22^{\text{sten}}$ Grades bezeichnen, setzen:

\weq{15}{b_x^6\, D(x,a) - a_x^6\, D(x,b) = f(x)\, \Phi(x,a,b).}

Denken wir uns in (15) für einen Augenblick ein neues Coordinatensystem eingeführt, in dem die Linien $a_x, b_x$ zwei Axen sind, von denen wir voraussetzen, dass sie sich nicht auf der Curve $f$ schneiden, so folgt aus der Vergleichung der rechten und linken Seite der identischen Gleichung (15), dass in $\Phi$ kein Glied vorkommen kann, was nicht entweder mit $a_x^6$ oder mit $b_x^6$ multiplicirt ist, und folglich hat $\Phi(x)$ die Form

\[
\Phi(x,a,b) = a_x^6\, \Phi_2(x) - b_x^6\, \Phi_1(x),
\]

worin $\Phi_1, \Phi_2$ Formen $16^{\text{ter}}$ Ordnung sind. Demnach erhält die identische Gleichung (15) die Form

\[
b_x^6\,[D(x,a) + f(x)\,\Phi_1(x)] = a_x^6\,[D(x,b) + f(x)\,\Phi_2(x)],
\]

und sie zeigt, dass $D(x,a) + f(x)\,\Phi_1(x)$ durch $a_x^6$ theilbar ist. Es existirt also eine Form $14^{\text{ten}}$ Grades $\chi(x,a,b)$, so dass

\weq{16}{\frac{D(x,a)}{a_x^6} = \chi(x,a,b) - f(x)\,\frac{\Phi_1(x,a,b)}{a_x^6}.}

Setzen wir hierin für die $a$ und $b$ irgend specielle, z. B. rationale Werthe $c, d$, und bezeichnen $\chi(x,c,d)$, was dann nur noch von $x$ und von den Coëfficienten der Form $f(x)$ abhängig ist, mit $\chi(x)$, so folgt:

\[
\frac{D(x,c)}{c_x^6} = \chi(x) - f(x)\,\frac{\Phi_1(x,c,d)}{c_x^6},
\]

und wenn wir dies von (16) subtrahiren und die Relation (15) für $b=c$ benutzen:

\weq{17}{\chi(x,a,b) - \chi(x) = f(x)\,\Psi(x),}

worin $\Psi(x)$ eine Function ist, die jedenfalls keinen anderen Nenner enthalten kann, als ein Product von Potenzen von $a_x$ und $c_x$. Da aber $f(x)$ nicht durch $a_x$ und $c_x$ theilbar ist, so muss $\Psi(x)$ eine ganze Function sein, und (16) ergiebt, wenn wir mit $\Theta(x)$ eine neue Form von $x$ (vom Grade 16) bezeichnen:

\weq{18}{\frac{D(x,a) - f(x)\,\Theta(x)}{a_x^6} = \chi(x).}
